\begin{lemma}[Uniform qualitative flatness]
Fix an integer $d\ge1$, put $n=d+1$ and identify $\R^n$ with the space-time
$\ell^2(\R^d\times\R)$. Let $b_d$ be the family flatness modulus of
\noderef{FlatnessModulus}. Then $b_d(s)\to0$ as $s\downarrow0$: for every
$\varepsilon>0$ there is $s_0$ with $0<s_0\le1/8$ and
\[
  b_d(s)\le\varepsilon\qquad(0<s\le s_0).
\]
\end{lemma}
\begin{proof}
Write $B_\rho$ for the open and $\overline B_\rho$ for the closed ball of radius
$\rho$ about the origin of the space-time $\ell^2(\R^d\times\R)$, put
$G=B_{3/4}$ and $D(z)=\operatorname{dist}(z,\partial G)$, and write
$e_b=(0,1)$, a unit vector. For $U$ in the normalized family $\mathcal U$ of
\noderef{NormalizedInfFamily}, $z\in G$ and $r>0$ let $\beta_U(z,r)$ be as in
\noderef{BetaFlatness}, so that by \noderef{FlatnessModulus}
\begin{equation}\label{eq:ff-bdef}
  b_d(s)=\sup\bigl\{\beta_U(z,r):U\in\mathcal U,\ z\in G,\ 0<r\le sD(z)\bigr\}.
\end{equation}
By \noderef{BetaFlatness} each $\beta_U(z,r)$ is the supremum of a subset of
$(0,1)$, hence $0\le\beta_U(z,r)\le1$; so the set in \eqref{eq:ff-bdef} is
bounded above by $1$, and therefore $0\le b_d(s)\le1$. What this paragraph
establishes is boundedness above, and nothing more: we do \emph{not} assert here
that either defining set is nonempty. Nonemptiness is needed only for the
approximation property of a supremum, it is used only in Step 2, and it is
derived there at the point of use, where the strict positivity of the suprema
concerned supplies it. Until then we use only the upper-bound property, which
boundedness alone licenses: $\beta_U(z,r)\le b_d(s)$ whenever $(U,z,r)$ is
admissible for $s$ as in \eqref{eq:ff-bdef}. We also use, as part of the
statement of \noderef{BetaFlatness} rather than as a consequence of it, that
every $\lambda\in(\beta_U(z,r),1)$ is good for $U$ at $(z,r)$, i.e.\ admits a
cylinder fit in the sense of \noderef{IsCylinderFit}.

Let $L\ge0$ be the common constant of \noderef{NormalizedFamilyLipschitz}, so
that every $U\in\mathcal U$ satisfies
\begin{equation}\label{eq:ff-lip}
  |U|\le L\ \text{ and }\
  |U(a)-U(b)|\le L|a-b|\qquad\bigl(a,b\in\overline B_{7/8}\bigr).
\end{equation}

Apart from the two suprema of \eqref{eq:ff-bdef} and of \noderef{BetaFlatness}
just discussed, every supremum written below is of the form
$\sup_{\overline B_\rho}|\cdot|$ over a closed ball of positive radius, hence
over a nonempty set, and in each case the function inside the absolute value is
bounded on that ball by \eqref{eq:ff-lip} --- together with the Cauchy--Schwarz
bound $|\langle Q,X\rangle|\le|Q|\,|X|$ in the instances where a fixed linear
functional also appears. Each of those suprema is therefore the supremum of a
nonempty set that is bounded above, so it is a genuine supremum, and we use this
below without further comment.

\emph{Step 0: every member of $\mathcal U$ is infinity-harmonic on $B_1$.}
Let $U\in\mathcal U$. By \noderef{NormalizedInfFamily} there is $u$,
infinity-harmonic on the unit ball of $\R^d$, with $U(z)=u(z_x)+z_t$ for every
space-time $z=(z_x,z_t)$. By \noderef{LinearLift} the function
$z\mapsto u(z_x)+z_t$ is infinity-harmonic on the preimage of that ball under
$z\mapsto z_x$, i.e.\ on $\{z:|z_x|<1\}$. Since $|z_x|\le|z|$, the ball $B_1$ is
contained in that set, so
\begin{equation}\label{eq:ff-harm}
  U\ \text{is infinity-harmonic on}\ B_1,\ \text{hence on}\ B_{7/8}.
\end{equation}
Here we used that infinity-harmonicity passes to an open subset: if
$\Omega'\subseteq\Omega$ are open and $v$ is infinity-harmonic on $\Omega$, then
for $z\in\Omega'$ and a $C^2$ function $\phi$ near $z$ such that $v-\phi$ has a
local maximum on $\Omega'$ at $z$, say $v-\phi\le(v-\phi)(z)$ on
$\mathcal N\cap\Omega'$ for a neighbourhood $\mathcal N$ of $z$, the set
$\mathcal N'=\mathcal N\cap\Omega'$ is again a neighbourhood of $z$ because
$\Omega'$ is open, and $\mathcal N'\cap\Omega=\mathcal N\cap\Omega'$; so
$v-\phi$ has a local maximum on $\Omega$ at $z$ and
\noderef{IsInfSubsolutionOn} gives $\Dinf\phi(z)\ge0$. The same argument with
minima and \noderef{IsInfSupersolutionOn} gives the other inequality, so $v$ is
infinity-harmonic on $\Omega'$ in the sense of \noderef{IsInfHarmonicOn}.

\emph{Step 1: a compactness claim.}
We shall use the following twice.

\textsc{Claim C.} \emph{Let $a$ be a point of a finite-dimensional real
inner-product space $\mathcal H$, let $\rho>0$, put
$K=\overline B(a,\rho)$, and let $v_i:\mathcal H\to\R$ satisfy
$|v_i|\le A$ on $K$ and $|v_i(w)-v_i(w')|\le\Lambda|w-w'|$ for $w,w'\in K$, with
constants $A,\Lambda\ge0$ independent of $i$. Then there are a strictly
increasing $\iota:\N\to\N$ and $v:\mathcal H\to\R$ with
$v_{\iota(k)}\to v$ uniformly on $K$.}

\emph{Proof of Claim C.} If $\mathcal H=\{0\}$ then $K=\{a\}$ is a single point,
the bounded sequence $(v_i(a))$ has a convergent subsequence by
Bolzano--Weierstrass (preamble), and the constant function with that limit
works. So assume $\mathcal H\ne\{0\}$. Fixing an orthonormal basis of
$\mathcal H$, the set of points with rational coordinates is countable and
dense, so its intersection with the open ball $B(a,\rho)$ is a countable set
$\{w_1,w_2,\dots\}$ dense in $B(a,\rho)$ and hence in its closure $K$. For each
$m$ the real sequence $(v_i(w_m))_i$ lies in $[-A,A]$; applying
Bolzano--Weierstrass (preamble) and extracting nested subsequences, and then
taking the diagonal, we obtain a strictly increasing $\iota$ such that
$(v_{\iota(k)}(w_m))_k$ converges for every $m$. We show
$(v_{\iota(k)})_k$ is uniformly Cauchy on $K$. Let $\varepsilon>0$ and put
$\sigma=\varepsilon/(3(\Lambda+1))>0$. The set $K$ is closed and bounded, hence
compact (preamble), so the cover of $K$ by the balls $B(w_m,\sigma)$ --- which
is a cover, by density --- has a finite subcover, indexed by
$m_1,\dots,m_p$. Choose $k_0$ with
$|v_{\iota(k)}(w_{m_l})-v_{\iota(k')}(w_{m_l})|<\varepsilon/3$ for all
$k,k'\ge k_0$ and all $l\le p$. Given $w\in K$ pick $l$ with
$|w-w_{m_l}|<\sigma$; then for $k,k'\ge k_0$
\[
  |v_{\iota(k)}(w)-v_{\iota(k')}(w)|
   \le\Lambda\sigma+\frac\varepsilon3+\Lambda\sigma\le\varepsilon,
\]
using the equi-Lipschitz bound twice and $\Lambda\sigma\le\varepsilon/3$. Hence
$(v_{\iota(k)})$ is uniformly Cauchy on $K$ and therefore converges uniformly on
$K$ to a function $v$ defined on $K$, which we extend to $\mathcal H$ by the
value $0$ off $K$. \hfill$\square$

\emph{Step 2: the contradiction hypothesis.}
Suppose the conclusion fails. Then there is $\varepsilon_*>0$ such that for
every $s_0\in(0,1/8]$ there is $s\in(0,s_0]$ with $b_d(s)>\varepsilon_*$. Put
$\delta=\varepsilon_*>0$ and apply this with $s_0=\min(1/8,1/i)$ for
$i=1,2,\dots$: we obtain $s_i\in(0,1/8]$ with
\begin{equation}\label{eq:ff-badb}
  s_i\le\frac1i\quad\text{and}\quad b_d(s_i)>\delta .
\end{equation}
In particular $s_i\to0$.

Both of the extractions we now perform use the approximation property of a
supremum, which is available only for a nonempty set; in both cases positivity
supplies nonemptiness. First, \noderef{FlatnessModulus} assigns $b_d(s_i)$ the
value $0$ when the set in \eqref{eq:ff-bdef} is empty, whereas
$b_d(s_i)>\delta>0$ by \eqref{eq:ff-badb}; hence that set is nonempty, and being
bounded above by $1$ it has a genuine supremum exceeding $\delta$. The
approximation property therefore applies and produces $U_i\in\mathcal U$,
$z_i\in G$ and $r_i$ with $0<r_i\le s_iD(z_i)$ and
\begin{equation}\label{eq:ff-beta}
  \beta_{U_i}(z_i,r_i)>\delta .
\end{equation}
Second, the same argument runs one level down. By \noderef{BetaFlatness} the
value $\beta_{U_i}(z_i,r_i)$ is $0$ when the set
$\{\lambda\in(0,1):\lambda\ \text{is not good for}\ U_i\ \text{at}\ (z_i,r_i)\}$
defining it is empty, whereas $\beta_{U_i}(z_i,r_i)>\delta>0$ by
\eqref{eq:ff-beta}; hence that set is nonempty too, and it is bounded above by
$1$. The approximation property applies to it as well and yields a parameter
\begin{equation}\label{eq:ff-lam}
  \lambda_i\in(\delta,1)\quad\text{that is \emph{not} good for }U_i
  \text{ at }(z_i,r_i).
\end{equation}
Since $D(z_i)\le3/4$ we get $r_i\le\tfrac34s_i\to0$.

\emph{Step 3: passing to a convergent subsequence.}
The points $z_i$ lie in the compact set $\overline B_{3/4}$, so after passing to
a subsequence $z_i\to z_*$ with $|z_*|\le3/4$. By \eqref{eq:ff-lip} the
functions $U_i$ are bounded by $L$ and $L$-Lipschitz on
$K=\overline B_{7/8}$, so Claim C (with $a=0$, $\rho=7/8$, $A=\Lambda=L$) gives,
after a further subsequence, a function $U_\circ$ with
\begin{equation}\label{eq:ff-unif}
  \eta_i:=\sup_{\overline B_{7/8}}|U_i-U_\circ|\longrightarrow0 .
\end{equation}
All the hypotheses recorded so far are preserved by passing to subsequences, and
we keep the same names. By \eqref{eq:ff-harm} each $U_i$ is infinity-harmonic on
the open set $B_{7/8}$, and \eqref{eq:ff-unif} gives uniform convergence on every
compact subset of $B_{7/8}$; so \noderef{InfHarmonicStability} applies and
\begin{equation}\label{eq:ff-limharm}
  U_\circ\ \text{is infinity-harmonic on}\ B_{7/8}.
\end{equation}
Note $z_*\in B_{7/8}$, since $|z_*|\le3/4<7/8$.

\emph{Step 4: the slow scales.}
Put $a_i=r_i+|z_i-z_*|+\eta_i$, so $a_i\to0$ by Steps 2 and 3. Passing to a
further subsequence we may assume $a_i\le4^{-i}$, and we set
\[
  R_i=2^{-i-4},
\]
so that $R_i\downarrow0$, $R_i\le1/16$, and
\begin{equation}\label{eq:ff-slow}
  \frac{a_i}{R_i}\le\frac{4^{-i}}{2^{-i-4}}=16\cdot2^{-i}\longrightarrow0 .
\end{equation}
Since $|z_*|\le3/4$ and $R_i\le1/16$ we have
\begin{equation}\label{eq:ff-inside}
  |z_*|+R_i\le\tfrac34+\tfrac1{16}<\tfrac78,
\end{equation}
so $z_*+R_iX\in B_{7/8}$ whenever $|X|\le1$.

\emph{Step 5: the blow-ups at $z_*$ are asymptotically linear.}
By \eqref{eq:ff-limharm} and $z_*\in B_{7/8}$ we may apply
\noderef{AffineBlowups} to the single function $U_\circ$ on $D=B_{7/8}$ at the
point $z_*$. Its second assertion, applied to the positive null sequence
$(R_i)$, gives a strictly increasing $\iota$ and a function $V$ such that the
blow-ups
\[
  \check U_k(X)=\frac{U_\circ(z_*+R_{\iota(k)}X)-U_\circ(z_*)}{R_{\iota(k)}}
\]
converge to $V$ uniformly on every compact set; its first assertion, applied to
the same data, gives a continuous linear functional $\Lambda_0$ with
$V=\Lambda_0$. By the Riesz representation of a continuous linear functional on
a finite-dimensional inner-product space there is $P_0$ with
$\Lambda_0(X)=\langle P_0,X\rangle$. Replace the whole family by the
$\iota$-subsequence and keep all names; \eqref{eq:ff-slow} and all earlier
properties persist. Thus
\begin{equation}\label{eq:ff-checkconv}
  \sup_{|X|\le1}\Bigl|
    \frac{U_\circ(z_*+R_iX)-U_\circ(z_*)}{R_i}-\langle P_0,X\rangle\Bigr|
  \longrightarrow0 .
\end{equation}
Define
\[
  \widehat U_i(X)=\frac{U_i(z_*+R_iX)-U_i(z_*)}{R_i}.
\]
For $|X|\le1$ both $z_*+R_iX$ and $z_*$ lie in $\overline B_{7/8}$ by
\eqref{eq:ff-inside}, so by \eqref{eq:ff-unif}
\[
  \bigl|\widehat U_i(X)-\check U_i(X)\bigr|
   \le\frac{2\eta_i}{R_i}\le\frac{2a_i}{R_i}\longrightarrow0
\]
by \eqref{eq:ff-slow}, where $\check U_i$ is as above with $R_i$ in place of
$R_{\iota(k)}$. Combining with \eqref{eq:ff-checkconv},
\begin{equation}\label{eq:ff-hatconv}
  \sup_{|X|\le1}\bigl|\widehat U_i(X)-\langle P_0,X\rangle\bigr|
  \longrightarrow0 .
\end{equation}

\emph{Step 6: the three-point estimate at the scales $R_i$.}
Each $\widehat U_i$ is infinity-harmonic on $B_1$. Indeed $U_i$ is
infinity-harmonic on $B_{7/8}$ by \eqref{eq:ff-harm};
\noderef{InfHarmonicRescale}, with $x_0=z_*$ and $A=R_i>0$, makes
$X\mapsto U_i(z_*+R_iX)$ infinity-harmonic on the preimage of $B_{7/8}$ under
$X\mapsto z_*+R_iX$, which contains $B_1$ by \eqref{eq:ff-inside}; and
\noderef{InfHarmonicNormalize}, with $c=U_i(z_*)$ and $M=R_i>0$, makes the
normalized difference $\widehat U_i$ infinity-harmonic on the same set. Also,
for $X,Y\in B_1$ the points $z_*+R_iX$ and $z_*+R_iY$ lie in
$\overline B_{7/8}$, so \eqref{eq:ff-lip} gives
\[
  |\widehat U_i(X)-\widehat U_i(Y)|
   =\frac{|U_i(z_*+R_iX)-U_i(z_*+R_iY)|}{R_i}\le L|X-Y| .
\]
Finally \eqref{eq:ff-hatconv} is exactly the hypothesis that
$\widehat U_i-\ell$ tends to $0$ uniformly on $B_1$ for the affine functions
$\ell(X)=\langle P_0,X\rangle+0$, the same one for every $i$. So
\noderef{ThreePointFlatness}, applied with this $L$, the functions
$\widehat U_i$, the slopes $P_0$ and the constants $0$, gives: for every
$\varepsilon>0$ there is $i_0$ with
\begin{equation}\label{eq:ff-threepoint}
  \bigl|\widehat U_i(X)+\widehat U_i(Y)
     -2\widehat U_i\bigl(\tfrac12(X+Y)\bigr)\bigr|\le\varepsilon|X-Y|
  \qquad\bigl(i\ge i_0,\ X,Y\in B_{1/4},\ X\ne Y\bigr).
\end{equation}

\emph{Step 7: the small-scale blow-ups.}
Since $r_i\to0$ we may assume $r_i\le1/16$ for every $i$, so that for $|H|\le2$
\[
  |z_i+r_iH|\le|z_i|+2r_i<\tfrac34+\tfrac18=\tfrac78 ,
\]
i.e.\ $z_i+r_iH\in\overline B_{7/8}$. Put
\[
  W_i(H)=\frac{U_i(z_i+r_iH)-U_i(z_i)}{r_i}.
\]
By \eqref{eq:ff-lip}, $|W_i(H)|\le L|H|\le2L$ and
$|W_i(H)-W_i(K)|\le L|H-K|$ for $H,K\in\overline B_2$. Claim C (with $a=0$,
$\rho=2$, $A=2L$, $\Lambda=L$) gives, after a further subsequence, a function
$W$ with
\begin{equation}\label{eq:ff-Wconv}
  \sup_{|H|\le2}|W_i-W|\longrightarrow0 .
\end{equation}
In particular $W(0)=\lim W_i(0)=0$, and, because $U_i(z)=u_i(z_x)+z_t$ for the
$u_i$ of \noderef{NormalizedInfFamily} and $e_b=(0,1)$ changes only the time
coordinate,
\begin{equation}\label{eq:ff-Web}
  W_i(e_b)=\frac{U_i(z_i+r_ie_b)-U_i(z_i)}{r_i}
    =\frac{\bigl(u_i((z_i)_x)+(z_i)_t+r_i\bigr)
            -\bigl(u_i((z_i)_x)+(z_i)_t\bigr)}{r_i}=1,
\end{equation}
so $W(e_b)=1$.

\emph{Step 8: $W$ has vanishing midpoint defect.}
Let $H,K\in\overline B_2$ with $H\ne K$ and put, for each $i$,
\[
  X_i=\frac{z_i-z_*+r_iH}{R_i},\qquad
  Y_i=\frac{z_i-z_*+r_iK}{R_i},
\]
so that $z_*+R_iX_i=z_i+r_iH$ and $z_*+R_iY_i=z_i+r_iK$, and
$\tfrac12(X_i+Y_i)=(z_i-z_*+r_i\tfrac12(H+K))/R_i$, whose image under
$X\mapsto z_*+R_iX$ is $z_i+r_i\tfrac12(H+K)$. Moreover
\[
  |X_i|\le\frac{|z_i-z_*|+2r_i}{R_i}\le\frac{3a_i}{R_i}\longrightarrow0
\]
by \eqref{eq:ff-slow}, and likewise for $Y_i$; so $X_i,Y_i\in B_{1/4}$ for all
large $i$. Also $X_i-Y_i=(r_i/R_i)(H-K)\ne0$, so $X_i\ne Y_i$.

Now the two midpoint defects are proportional. Since
$\widehat U_i(X)=(U_i(z_*+R_iX)-U_i(z_*))/R_i$ and the three coefficients
$1,1,-2$ sum to zero, the constants $U_i(z_*)$ cancel and
\[
  \widehat U_i(X_i)+\widehat U_i(Y_i)
    -2\widehat U_i\bigl(\tfrac12(X_i+Y_i)\bigr)
  =\frac{U_i(z_i+r_iH)+U_i(z_i+r_iK)
          -2U_i\bigl(z_i+r_i\tfrac12(H+K)\bigr)}{R_i};
\]
the same cancellation of $U_i(z_i)$ in the definition of $W_i$ gives
\[
  W_i(H)+W_i(K)-2W_i\bigl(\tfrac12(H+K)\bigr)
  =\frac{U_i(z_i+r_iH)+U_i(z_i+r_iK)
          -2U_i\bigl(z_i+r_i\tfrac12(H+K)\bigr)}{r_i}.
\]
Hence the left sides differ by the factor $R_i/r_i$. Let $\varepsilon>0$ and let
$i_0$ be as in \eqref{eq:ff-threepoint}. For $i\ge i_0$ large enough that
$X_i,Y_i\in B_{1/4}$,
\[
  \bigl|W_i(H)+W_i(K)-2W_i\bigl(\tfrac12(H+K)\bigr)\bigr|
   =\frac{R_i}{r_i}\bigl|\widehat U_i(X_i)+\widehat U_i(Y_i)
      -2\widehat U_i\bigl(\tfrac12(X_i+Y_i)\bigr)\bigr|
   \le\frac{R_i}{r_i}\varepsilon|X_i-Y_i|=\varepsilon|H-K| ,
\]
using $|X_i-Y_i|=(r_i/R_i)|H-K|$. Letting $i\to\infty$ and using
\eqref{eq:ff-Wconv} at the three points $H$, $K$, $\tfrac12(H+K)$ of
$\overline B_2$,
\[
  \bigl|W(H)+W(K)-2W\bigl(\tfrac12(H+K)\bigr)\bigr|\le\varepsilon|H-K| .
\]
As $\varepsilon>0$ was arbitrary, and the case $H=K$ is trivial,
\begin{equation}\label{eq:ff-midpoint}
  W(H)+W(K)=2W\bigl(\tfrac12(H+K)\bigr)\qquad\bigl(H,K\in\overline B_2\bigr).
\end{equation}

\emph{Step 9: $W$ is linear.}
$W$ is $L$-Lipschitz on $\overline B_2$, being a uniform limit of $L$-Lipschitz
functions, and $W(0)=0$. Taking $K=0$ in \eqref{eq:ff-midpoint} and using
$W(0)=0$,
\begin{equation}\label{eq:ff-half}
  W(H)=2W(H/2)\qquad\bigl(H\in\overline B_2\bigr).
\end{equation}
For $H,K\in\overline B_1$ we have $H+K\in\overline B_2$ and
$\tfrac12(H+K)\in\overline B_1$, so \eqref{eq:ff-midpoint} and
\eqref{eq:ff-half} give
\begin{equation}\label{eq:ff-add}
  W(H)+W(K)=2W\bigl(\tfrac12(H+K)\bigr)=W(H+K)
  \qquad\bigl(H,K\in\overline B_1\bigr).
\end{equation}
In particular $W(-H)=-W(H)$ for $H\in\overline B_1$ (take $K=-H$ and use
$W(0)=0$). Next,
\begin{equation}\label{eq:ff-dyadic}
  W(qH)=qW(H)\qquad\bigl(q\in[0,1]\ \text{a dyadic rational},\
  H\in\overline B_1\bigr).
\end{equation}
We prove \eqref{eq:ff-dyadic} by subdividing, not by doubling, and subdivision
is the only direction the identities support: \eqref{eq:ff-add} is stated for
$H,K\in\overline B_1$, so taking $K=H$ in it gives $W(2H)=2W(H)$ \emph{once},
but applying it a second time would require $2H\in\overline B_1$, which fails as
soon as $|H|>1/2$. Fix $H\in\overline B_1$ and write $q=p2^{-m}$ with integers
$m\ge0$ and $0\le p\le2^m$. Halving first: $H\in\overline B_1\subseteq
\overline B_2$ and $\overline B_2$ is closed under $X\mapsto X/2$, so
\eqref{eq:ff-half} may be applied $m$ times in succession, each time to a point
of $\overline B_2$, and it gives
\[
  W\bigl(2^{-m}H\bigr)=2^{-m}W(H).
\]
Adding second: form the partial sums $k2^{-m}H$ for $k=0,1,\dots,p$. For
$0\le k<p$ we have $\bigl|k2^{-m}H\bigr|=k2^{-m}|H|\le q|H|\le|H|\le1$ and
$\bigl|2^{-m}H\bigr|\le|H|\le1$, so the partial sum and the increment both lie
in $\overline B_1$ and \eqref{eq:ff-add} applies at each of the $p$ steps. Hence
\[
  W\bigl(p2^{-m}H\bigr)=pW\bigl(2^{-m}H\bigr)=p2^{-m}W(H)=qW(H),
\]
which is \eqref{eq:ff-dyadic}. (This is exactly the bookkeeping used for the
basis expansion in the next paragraph, performed one step earlier.) Since the
dyadic rationals are dense in $[0,1]$, the continuity of $W$ extends
\eqref{eq:ff-dyadic} to
\begin{equation}\label{eq:ff-hom}
  W(tH)=tW(H)\qquad\bigl(t\in[-1,1],\ H\in\overline B_1\bigr),
\end{equation}
the negative values following from oddness.

Let $e_1,\dots,e_n$ be an orthonormal basis of the space-time, and set
$P=\sum_{j}W(e_j)e_j$. Let $H\in\overline B_1$ and write
$H=\sum_ja_je_j$ with $a_j=\langle H,e_j\rangle$. Every partial sum
$\sum_{j\le k}a_je_j$ has norm at most $|H|\le1$, and each $a_ke_k$ has norm
$|a_k|\le1$, so \eqref{eq:ff-add} applies repeatedly and gives
$W(H)=\sum_jW(a_je_j)$; and $W(a_je_j)=a_jW(e_j)$ by \eqref{eq:ff-hom}.
Therefore
\[
  W(H)=\sum_ja_jW(e_j)=\langle P,H\rangle\qquad\bigl(H\in\overline B_1\bigr),
\]
and then \eqref{eq:ff-half} extends this to $\overline B_2$:
$W(H)=2W(H/2)=2\langle P,H/2\rangle=\langle P,H\rangle$. By
\eqref{eq:ff-Web}, $\langle P,e_b\rangle=W(e_b)=1$, so in particular
\begin{equation}\label{eq:ff-P}
  W(H)=\langle P,H\rangle\ \bigl(H\in\overline B_2\bigr),\qquad P\ne0 .
\end{equation}

\emph{Step 10: the contradiction.}
By \eqref{eq:ff-Wconv} and \eqref{eq:ff-P} there is $i$ with
\begin{equation}\label{eq:ff-final}
  \sup_{|H|\le2}\bigl|U_i(z_i+r_iH)-U_i(z_i)-r_i\langle P,H\rangle\bigr|
   =r_i\sup_{|H|\le2}\bigl|W_i(H)-\langle P,H\rangle\bigr|<\delta r_i ,
\end{equation}
and we fix such an $i$, large enough for all the finitely many "for all large
$i$" requirements above. We claim $P$ realizes the cylinder fit of parameter
$\lambda_i$ for $U_i$ at $(z_i,r_i)$ in the sense of \noderef{IsCylinderFit}.
First $P\ne0$ by \eqref{eq:ff-P}. Next let $Z$ belong to the cylinder
$Z_{\lambda_i}(z_i,r_i;P)$ of \noderef{FitCylinder} and write $Z=z_i+h$,
$N=P/|P|$. By that definition $|\langle N,h\rangle|\le r_i$ and
$|h-\langle N,h\rangle N|\le\sqrt{\lambda_i}\,r_i$, and these two components of
$h$ are orthogonal, so
\[
  |h|^2=|\langle N,h\rangle|^2+|h-\langle N,h\rangle N|^2
   \le r_i^2+\lambda_ir_i^2\le2r_i^2 ,
\]
using $\lambda_i<1$; hence $|h|\le\sqrt2\,r_i\le2r_i$. Writing $H=h/r_i$, so
$|H|\le2$, \eqref{eq:ff-final} gives
\[
  \bigl|U_i(Z)-U_i(z_i)-\langle P,Z-z_i\rangle\bigr|
   =\bigl|U_i(z_i+r_iH)-U_i(z_i)-r_i\langle P,H\rangle\bigr|
   <\delta r_i\le\lambda_ir_i ,
\]
the last step by $\delta<\lambda_i$ from \eqref{eq:ff-lam}. So $\lambda_i$ is
good for $U_i$ at $(z_i,r_i)$, contradicting \eqref{eq:ff-lam}. Therefore the
conclusion cannot fail.
\end{proof}
